% Options for packages loaded elsewhere
\PassOptionsToPackage{unicode}{hyperref}
\PassOptionsToPackage{hyphens}{url}
\documentclass[
]{article}
\usepackage{xcolor}
\usepackage[margin=1in]{geometry}
\usepackage{amsmath,amssymb}
\setcounter{secnumdepth}{5}
\usepackage{iftex}
\ifPDFTeX
  \usepackage[T1]{fontenc}
  \usepackage[utf8]{inputenc}
  \usepackage{textcomp} % provide euro and other symbols
\else % if luatex or xetex
  \usepackage{unicode-math} % this also loads fontspec
  \defaultfontfeatures{Scale=MatchLowercase}
  \defaultfontfeatures[\rmfamily]{Ligatures=TeX,Scale=1}
\fi
\usepackage{lmodern}
\ifPDFTeX\else
  % xetex/luatex font selection
\fi
% Use upquote if available, for straight quotes in verbatim environments
\IfFileExists{upquote.sty}{\usepackage{upquote}}{}
\IfFileExists{microtype.sty}{% use microtype if available
  \usepackage[]{microtype}
  \UseMicrotypeSet[protrusion]{basicmath} % disable protrusion for tt fonts
}{}
\makeatletter
\@ifundefined{KOMAClassName}{% if non-KOMA class
  \IfFileExists{parskip.sty}{%
    \usepackage{parskip}
  }{% else
    \setlength{\parindent}{0pt}
    \setlength{\parskip}{6pt plus 2pt minus 1pt}}
}{% if KOMA class
  \KOMAoptions{parskip=half}}
\makeatother
\usepackage{color}
\usepackage{fancyvrb}
\newcommand{\VerbBar}{|}
\newcommand{\VERB}{\Verb[commandchars=\\\{\}]}
\DefineVerbatimEnvironment{Highlighting}{Verbatim}{commandchars=\\\{\}}
% Add ',fontsize=\small' for more characters per line
\usepackage{framed}
\definecolor{shadecolor}{RGB}{248,248,248}
\newenvironment{Shaded}{\begin{snugshade}}{\end{snugshade}}
\newcommand{\AlertTok}[1]{\textcolor[rgb]{0.94,0.16,0.16}{#1}}
\newcommand{\AnnotationTok}[1]{\textcolor[rgb]{0.56,0.35,0.01}{\textbf{\textit{#1}}}}
\newcommand{\AttributeTok}[1]{\textcolor[rgb]{0.13,0.29,0.53}{#1}}
\newcommand{\BaseNTok}[1]{\textcolor[rgb]{0.00,0.00,0.81}{#1}}
\newcommand{\BuiltInTok}[1]{#1}
\newcommand{\CharTok}[1]{\textcolor[rgb]{0.31,0.60,0.02}{#1}}
\newcommand{\CommentTok}[1]{\textcolor[rgb]{0.56,0.35,0.01}{\textit{#1}}}
\newcommand{\CommentVarTok}[1]{\textcolor[rgb]{0.56,0.35,0.01}{\textbf{\textit{#1}}}}
\newcommand{\ConstantTok}[1]{\textcolor[rgb]{0.56,0.35,0.01}{#1}}
\newcommand{\ControlFlowTok}[1]{\textcolor[rgb]{0.13,0.29,0.53}{\textbf{#1}}}
\newcommand{\DataTypeTok}[1]{\textcolor[rgb]{0.13,0.29,0.53}{#1}}
\newcommand{\DecValTok}[1]{\textcolor[rgb]{0.00,0.00,0.81}{#1}}
\newcommand{\DocumentationTok}[1]{\textcolor[rgb]{0.56,0.35,0.01}{\textbf{\textit{#1}}}}
\newcommand{\ErrorTok}[1]{\textcolor[rgb]{0.64,0.00,0.00}{\textbf{#1}}}
\newcommand{\ExtensionTok}[1]{#1}
\newcommand{\FloatTok}[1]{\textcolor[rgb]{0.00,0.00,0.81}{#1}}
\newcommand{\FunctionTok}[1]{\textcolor[rgb]{0.13,0.29,0.53}{\textbf{#1}}}
\newcommand{\ImportTok}[1]{#1}
\newcommand{\InformationTok}[1]{\textcolor[rgb]{0.56,0.35,0.01}{\textbf{\textit{#1}}}}
\newcommand{\KeywordTok}[1]{\textcolor[rgb]{0.13,0.29,0.53}{\textbf{#1}}}
\newcommand{\NormalTok}[1]{#1}
\newcommand{\OperatorTok}[1]{\textcolor[rgb]{0.81,0.36,0.00}{\textbf{#1}}}
\newcommand{\OtherTok}[1]{\textcolor[rgb]{0.56,0.35,0.01}{#1}}
\newcommand{\PreprocessorTok}[1]{\textcolor[rgb]{0.56,0.35,0.01}{\textit{#1}}}
\newcommand{\RegionMarkerTok}[1]{#1}
\newcommand{\SpecialCharTok}[1]{\textcolor[rgb]{0.81,0.36,0.00}{\textbf{#1}}}
\newcommand{\SpecialStringTok}[1]{\textcolor[rgb]{0.31,0.60,0.02}{#1}}
\newcommand{\StringTok}[1]{\textcolor[rgb]{0.31,0.60,0.02}{#1}}
\newcommand{\VariableTok}[1]{\textcolor[rgb]{0.00,0.00,0.00}{#1}}
\newcommand{\VerbatimStringTok}[1]{\textcolor[rgb]{0.31,0.60,0.02}{#1}}
\newcommand{\WarningTok}[1]{\textcolor[rgb]{0.56,0.35,0.01}{\textbf{\textit{#1}}}}
\usepackage{graphicx}
\makeatletter
\newsavebox\pandoc@box
\newcommand*\pandocbounded[1]{% scales image to fit in text height/width
  \sbox\pandoc@box{#1}%
  \Gscale@div\@tempa{\textheight}{\dimexpr\ht\pandoc@box+\dp\pandoc@box\relax}%
  \Gscale@div\@tempb{\linewidth}{\wd\pandoc@box}%
  \ifdim\@tempb\p@<\@tempa\p@\let\@tempa\@tempb\fi% select the smaller of both
  \ifdim\@tempa\p@<\p@\scalebox{\@tempa}{\usebox\pandoc@box}%
  \else\usebox{\pandoc@box}%
  \fi%
}
% Set default figure placement to htbp
\def\fps@figure{htbp}
\makeatother
\setlength{\emergencystretch}{3em} % prevent overfull lines
\providecommand{\tightlist}{%
  \setlength{\itemsep}{0pt}\setlength{\parskip}{0pt}}
\usepackage{bookmark}
\IfFileExists{xurl.sty}{\usepackage{xurl}}{} % add URL line breaks if available
\urlstyle{same}
\hypersetup{
  pdftitle={Coding challenge 06: Functions and Iterations},
  pdfauthor={Dulshika Dassanayake},
  hidelinks,
  pdfcreator={LaTeX via pandoc}}

\title{Coding challenge 06: Functions and Iterations}
\author{Dulshika Dassanayake}
\date{2026-03-23}

\begin{document}
\maketitle

{
\setcounter{tocdepth}{2}
\tableofcontents
}
\section{Regarding reproducibility, what is the main point of writing
your own functions and
iterations?}\label{regarding-reproducibility-what-is-the-main-point-of-writing-your-own-functions-and-iterations}

Writing own functions and iterations helps ensure reproducibility by
making the analysis systematic and repeatable. Functions allow the same
set of instructions to be applied consistently each time. Iterations
automate repeating the same process across multiple datasets or
variables. This reduces human error and ensures the results can be
reproduced exactly.

\section{Conceptual:}\label{conceptual}

A function in R is written using the function() command. The inputs
(arguments) are placed inside the parentheses, and the code that
performs the calculation is written inside curly brackets \{ \}. The
function processes the inputs and returns the result using return() or
by outputting the final value.

A for loop in R is written using the syntax for (i in sequence). The
loop runs the same block of code repeatedly for each element in the
sequence. The code inside the curly brackets executes during each
iteration, and results are typically stored in objects such as vectors
or data frames.

\section{Read data}\label{read-data}

\begin{Shaded}
\begin{Highlighting}[]
\FunctionTok{library}\NormalTok{(ggplot2)}
\FunctionTok{library}\NormalTok{(drc) }
\end{Highlighting}
\end{Shaded}

\begin{verbatim}
## Loading required package: MASS
\end{verbatim}

\begin{verbatim}
## 
## 'drc' has been loaded.
\end{verbatim}

\begin{verbatim}
## Please cite R and 'drc' if used for a publication,
\end{verbatim}

\begin{verbatim}
## for references type 'citation()' and 'citation('drc')'.
\end{verbatim}

\begin{verbatim}
## 
## Attaching package: 'drc'
\end{verbatim}

\begin{verbatim}
## The following objects are masked from 'package:stats':
## 
##     gaussian, getInitial
\end{verbatim}

\begin{Shaded}
\begin{Highlighting}[]
\FunctionTok{library}\NormalTok{(tidyverse)}
\end{Highlighting}
\end{Shaded}

\begin{verbatim}
## -- Attaching core tidyverse packages ------------------------ tidyverse 2.0.0 --
## v dplyr     1.1.4     v readr     2.1.6
## v forcats   1.0.1     v stringr   1.6.0
## v lubridate 1.9.4     v tibble    3.3.1
## v purrr     1.2.1     v tidyr     1.3.2
\end{verbatim}

\begin{verbatim}
## -- Conflicts ------------------------------------------ tidyverse_conflicts() --
## x dplyr::filter() masks stats::filter()
## x dplyr::lag()    masks stats::lag()
## x dplyr::select() masks MASS::select()
## i Use the conflicted package (<http://conflicted.r-lib.org/>) to force all conflicts to become errors
\end{verbatim}

\begin{Shaded}
\begin{Highlighting}[]
\NormalTok{citydata}\OtherTok{=}\FunctionTok{read.csv}\NormalTok{ (}\StringTok{"Cities.csv"}\NormalTok{, }\AttributeTok{na.strings =} \StringTok{"na"}\NormalTok{)}
\FunctionTok{head}\NormalTok{(citydata)}
\end{Highlighting}
\end{Shaded}

\begin{verbatim}
##          city  city_ascii state_id state_name county_fips county_name     lat
## 1    New York    New York       NY   New York       36081      Queens 40.6943
## 2 Los Angeles Los Angeles       CA California        6037 Los Angeles 34.1141
## 3     Chicago     Chicago       IL   Illinois       17031        Cook 41.8375
## 4       Miami       Miami       FL    Florida       12086  Miami-Dade 25.7840
## 5     Houston     Houston       TX      Texas       48201      Harris 29.7860
## 6      Dallas      Dallas       TX      Texas       48113      Dallas 32.7935
##        long population density
## 1  -73.9249   18832416 10943.7
## 2 -118.4068   11885717  3165.8
## 3  -87.6866    8489066  4590.3
## 4  -80.2101    6113982  4791.1
## 5  -95.3885    6046392  1386.5
## 6  -96.7667    5843632  1477.2
\end{verbatim}

\section{Function writing}\label{function-writing}

\begin{Shaded}
\begin{Highlighting}[]
\NormalTok{distance }\OtherTok{\textless{}{-}} \ControlFlowTok{function}\NormalTok{(lat1, lon1, lat2, lon2)\{}
  \CommentTok{\#convert to radians}
\NormalTok{  rad.lat1 }\OtherTok{\textless{}{-}}\NormalTok{ lat1 }\SpecialCharTok{*}\NormalTok{ pi}\SpecialCharTok{/}\DecValTok{180}
\NormalTok{  rad.lon1 }\OtherTok{\textless{}{-}}\NormalTok{ lon1 }\SpecialCharTok{*}\NormalTok{ pi}\SpecialCharTok{/}\DecValTok{180}
\NormalTok{  rad.lat2 }\OtherTok{\textless{}{-}}\NormalTok{ lat2 }\SpecialCharTok{*}\NormalTok{ pi}\SpecialCharTok{/}\DecValTok{180}
\NormalTok{  rad.lon2 }\OtherTok{\textless{}{-}}\NormalTok{ lon2 }\SpecialCharTok{*}\NormalTok{ pi}\SpecialCharTok{/}\DecValTok{180}
  
  \CommentTok{\#Haversine formula}
\NormalTok{  delta\_lat }\OtherTok{\textless{}{-}}\NormalTok{ rad.lat2 }\SpecialCharTok{{-}}\NormalTok{ rad.lat1}
\NormalTok{  delta\_lon }\OtherTok{\textless{}{-}}\NormalTok{ rad.lon2 }\SpecialCharTok{{-}}\NormalTok{ rad.lon1}
\NormalTok{  a }\OtherTok{\textless{}{-}} \FunctionTok{sin}\NormalTok{(delta\_lat }\SpecialCharTok{/} \DecValTok{2}\NormalTok{)}\SpecialCharTok{\^{}}\DecValTok{2} \SpecialCharTok{+} \FunctionTok{cos}\NormalTok{(rad.lat1) }\SpecialCharTok{*} \FunctionTok{cos}\NormalTok{(rad.lat2) }\SpecialCharTok{*} \FunctionTok{sin}\NormalTok{(delta\_lon }\SpecialCharTok{/} \DecValTok{2}\NormalTok{)}\SpecialCharTok{\^{}}\DecValTok{2}
\NormalTok{  c }\OtherTok{\textless{}{-}} \DecValTok{2} \SpecialCharTok{*} \FunctionTok{asin}\NormalTok{(}\FunctionTok{sqrt}\NormalTok{(a))}
  
  \CommentTok{\#Earth\textquotesingle{}s radius in km}
\NormalTok{  earth\_radius }\OtherTok{\textless{}{-}} \DecValTok{6378137}
  
  \CommentTok{\#Calculate the distance}
\NormalTok{  distance\_km }\OtherTok{\textless{}{-}}\NormalTok{ (earth\_radius }\SpecialCharTok{*}\NormalTok{ c)}\SpecialCharTok{/}\DecValTok{1000}
  \FunctionTok{return}\NormalTok{(distance\_km)}
\NormalTok{\}}
\end{Highlighting}
\end{Shaded}

\section{calculate the distance between Auburn, AL and New York
City.}\label{calculate-the-distance-between-auburn-al-and-new-york-city.}

\begin{Shaded}
\begin{Highlighting}[]
\NormalTok{Au\_NY }\OtherTok{\textless{}{-}} \ControlFlowTok{function}\NormalTok{(start, destination)}
\NormalTok{\{}
\NormalTok{  cord\_extract }\OtherTok{\textless{}{-}}\NormalTok{ citydata }\SpecialCharTok{\%\textgreater{}\%} 
  \FunctionTok{select}\NormalTok{(city, lat, long) }\SpecialCharTok{\%\textgreater{}\%}
  \FunctionTok{filter}\NormalTok{(city }\SpecialCharTok{==}\NormalTok{ start }\SpecialCharTok{|}\NormalTok{ city }\SpecialCharTok{==}\NormalTok{ destination) }\SpecialCharTok{\%\textgreater{}\%}
  \FunctionTok{group\_by}\NormalTok{(city)}
\FunctionTok{head}\NormalTok{(cord\_extract)}
\NormalTok{cords }\OtherTok{\textless{}{-}} \FunctionTok{list}\NormalTok{(cord\_extract}\SpecialCharTok{$}\NormalTok{lat, cord\_extract}\SpecialCharTok{$}\NormalTok{long)}
\NormalTok{latitude\_1 }\OtherTok{\textless{}{-}}\NormalTok{ cords[[}\DecValTok{1}\NormalTok{]][}\DecValTok{1}\NormalTok{]}
\NormalTok{latitude\_2 }\OtherTok{\textless{}{-}}\NormalTok{ cords[[}\DecValTok{1}\NormalTok{]][}\DecValTok{2}\NormalTok{]}
\NormalTok{longitude\_1 }\OtherTok{\textless{}{-}}\NormalTok{ cords[[}\DecValTok{2}\NormalTok{]][}\DecValTok{1}\NormalTok{]}
\NormalTok{longitude\_2 }\OtherTok{\textless{}{-}}\NormalTok{ cords[[}\DecValTok{2}\NormalTok{]][}\DecValTok{2}\NormalTok{]}
\FunctionTok{return}\NormalTok{(}\FunctionTok{distance}\NormalTok{(latitude\_1, longitude\_1, latitude\_2, longitude\_2))}
\NormalTok{\}}

\FunctionTok{Au\_NY}\NormalTok{ (}\StringTok{"New York"}\NormalTok{, }\StringTok{"Auburn"}\NormalTok{)}
\end{Highlighting}
\end{Shaded}

\begin{verbatim}
## [1] 1367.854
\end{verbatim}

\section{for loop to calculate the distance between Auburn, AL and every
other city in
Cities.csv}\label{for-loop-to-calculate-the-distance-between-auburn-al-and-every-other-city-in-cities.csv}

\begin{Shaded}
\begin{Highlighting}[]
\NormalTok{Au\_to\_cities }\OtherTok{\textless{}{-}} \FunctionTok{unique}\NormalTok{(citydata}\SpecialCharTok{$}\NormalTok{city)}
\ControlFlowTok{for}\NormalTok{ (i }\ControlFlowTok{in} \FunctionTok{seq\_along}\NormalTok{(Au\_to\_cities)) \{}
\NormalTok{  results }\OtherTok{\textless{}{-}} \FunctionTok{Au\_NY}\NormalTok{(Au\_to\_cities[[i]], }\StringTok{"Auburn"}\NormalTok{ )}
  \FunctionTok{print}\NormalTok{(results)}
\NormalTok{\}}
\end{Highlighting}
\end{Shaded}

\begin{verbatim}
## [1] 1367.854
## [1] 3051.838
## [1] 1045.521
## [1] 916.4138
## [1] 993.0298
## [1] 1056.022
## [1] 1239.973
## [1] 162.5121
## [1] 1036.99
## [1] 1665.699
## [1] 2476.255
## [1] 1108.229
## [1] 3507.959
## [1] 3388.366
## [1] 2951.382
## [1] 1530.2
## [1] 591.1181
## [1] 1363.207
## [1] 1909.79
## [1] 1380.138
## [1] 2961.12
## [1] 2752.814
## [1] 1092.259
## [1] 796.7541
## [1] 3479.538
## [1] 1290.549
## [1] 3301.992
## [1] 1191.666
## [1] 608.2035
## [1] 2504.631
## [1] 3337.278
## [1] 800.1452
## [1] 1001.088
## [1] 732.5906
## [1] 1371.163
## [1] 1091.897
## [1] 1043.273
## [1] 851.3423
## [1] 1382.372
## [1] NA
\end{verbatim}

\section{Modify loop}\label{modify-loop}

\begin{Shaded}
\begin{Highlighting}[]
\NormalTok{all\_distances }\OtherTok{\textless{}{-}} \ConstantTok{NULL}
\NormalTok{Au\_to\_cities }\OtherTok{\textless{}{-}} \FunctionTok{unique}\NormalTok{(citydata}\SpecialCharTok{$}\NormalTok{city)}
\ControlFlowTok{for}\NormalTok{ (i }\ControlFlowTok{in} \FunctionTok{seq\_along}\NormalTok{(Au\_to\_cities)) \{}
\NormalTok{  results }\OtherTok{\textless{}{-}} \FunctionTok{Au\_NY}\NormalTok{(Au\_to\_cities[[i]], }\StringTok{"Auburn"}\NormalTok{ )}
\NormalTok{  distance\_i }\OtherTok{\textless{}{-}} \FunctionTok{data.frame}\NormalTok{(Au\_to\_cities[[i]], }\StringTok{"Auburn"}\NormalTok{, results)}
  \FunctionTok{colnames}\NormalTok{(distance\_i) }\OtherTok{\textless{}{-}} \FunctionTok{c}\NormalTok{(}\StringTok{"Destination"}\NormalTok{, }\StringTok{"Origin"}\NormalTok{, }\StringTok{"Distance (Km)"}\NormalTok{)}
\NormalTok{  all\_distances }\OtherTok{\textless{}{-}} \FunctionTok{rbind.data.frame}\NormalTok{(all\_distances, distance\_i)}
\NormalTok{\}}

\FunctionTok{head}\NormalTok{(all\_distances)}
\end{Highlighting}
\end{Shaded}

\begin{verbatim}
##   Destination Origin Distance (Km)
## 1    New York Auburn     1367.8540
## 2 Los Angeles Auburn     3051.8382
## 3     Chicago Auburn     1045.5213
## 4       Miami Auburn      916.4138
## 5     Houston Auburn      993.0298
## 6      Dallas Auburn     1056.0217
\end{verbatim}

\section{GitHub}\label{github}

Github my Github repository:

\end{document}
